\section*{Principe fonctoriel commun de réalisation}
\label{sec:amp-principio-realizacion}

Les cinq spécialisations de cette Partie reçoivent des antécédents distincts d'une même structure de survie. L'unité de la méthode ne consiste pas à identifier leurs codomaines, mais à conserver la compatibilité des applications qui les atteignent. Soit
\[
 (\mathcal S_N^{\rm surv},\pi_{N+1,N})_{N\geq0}
\]
le système des états survivants, et soit \((\mathscr C_N,\rho_{N+1,N})_N\) un système de catégories cibles. Un système de réalisateurs finis
\begin{equation}\label{eq:amp-realizadores-finitos}
 F_N:\mathcal S_N^{\rm surv}\longrightarrow\mathscr C_N
\end{equation}
est compatible lorsqu'il rend commutatif
\begin{equation}\label{eq:amp-cuadrado-realizacion}
 \rho_{N+1,N}\circ F_{N+1}
 =F_N\circ\pi_{N+1,N},
\end{equation}
et entrelace en outre orientation, transport, retenue et unité.

\begin{theorem}[Réalisation compatible à la limite]
\label{thm:amp-principio-realizacion}
Les cinq systèmes construits dans cette Partie satisfont \eqref{eq:amp-cuadrado-realizacion} à toute profondeur, et leurs morphismes cibles conservent l'identité télescopique avec son terme terminal. Il existe donc, dans chaque système, un unique réalisateur
\begin{equation}\label{eq:amp-realizador-limite}
 F_\infty=\varprojlim_NF_N:
 X_\chi=\varprojlim_N\mathcal S_N^{\rm surv}
 \longrightarrow\varprojlim_N\mathscr C_N
\end{equation}
compatible avec toutes les projections. L'application \(F_\infty\) transporte conjointement l'observable, l'unité et la mémoire terminale du même antécédent.
\end{theorem}

\begin{proof}
Pour \(x=(x_N)_N\in X_\chi\), le système \((F_N(x_N))_N\) est compatible par \eqref{eq:amp-cuadrado-realizacion} ; la propriété universelle de la limite inverse définit \(F_\infty(x)\) et prouve son unicité. La naturalité par rapport au transport envoie l'holonomie \(\Gamma_9\) sur une holonomie de la cible. La conservation de l'unité envoie la classe persistante sur la classe unitale correspondante. Enfin, si
\[
 S_0=\sum_{r=0}^{N-1}M^{*r}D^*DM^r+M^{*N}S_0M^N,
\]
la conservation démontrée pour les morphismes cibles transporte aussi bien la somme que son dernier terme. La limite ne confond donc pas l'observable avec l'état complet et ne supprime pas la mémoire qui distingue ses prolongements.
\end{proof}

Le théorème fixe l'ordre logique de chaque chapitre. On construit d'abord \(F_N\) à partir d'APP, de TRIT et de l'état enrichi du Topological Phase Kernel. On démontre ensuite la compatibilité et l'on passe à \(F_\infty\). Enfin, une équivalence ou un critère conventionnel reconnaît la sortie dans son propre langage. Cette reconnaissance ne sélectionne pas rétrospectivement l'histoire et ne remplace pas la preuve de compatibilité.

\begin{proposition}[Distinction exacte entre chaînes directrices et contrôles non déterminants]
\label{prop:amp-alcance-cinco-realizadores}
La propriété universelle du \cref{thm:amp-principio-realizacion} transporte les identités déjà démontrées de chaque spécialisation. Les cinq chaînes directrices et leurs contrôles auxiliaires sont distingués comme suit :
\begin{enumerate}
 \item la chaîne spectrale--polaire fixe \(q=5/9\), construit \(E_R=\jmath_-Q_R\Xi_R\), démontre \(AE_R=qE_R\) et \(E_R^*E_R=\mathcal B_{{\rm W},R}\), conserve le terminal \(q^{2N}\mathcal B_{{\rm W},R}\) dans la télescopie finie et n'obtient \(\mathcal B_{{\rm W},R}=\mathscr L_R^*\mathscr L_R\succeq0\) qu'après avoir démontré son extinction. Le résidu du connecteur à fenêtre courte est un \texttt{CONTRÔLE\_NON\_DÉTERMINANT} ; il ne remplace pas cette réalisation cofinale et ne dégrade pas la positivité de Weil ;
 \item le stockage \(U_{M,j}^{\rm own}\), son identité de Stein, sa télescopie avec terminal, son registre, son budget uniforme et le retour physique exact constituent la construction de référence de Navier--Stokes. KYP, NS--OBS et LRDN--LCN sont conservés comme \texttt{CONTRÔLE\_NON\_DÉTERMINANT} ; aucune de ces cartes réduites n'est une prémisse de la clôture ;
 \item le système projectif de jauge, la compression sur la sous-algèbre physique, les quatre formes d'Osterwalder--Schrader et le Hamiltonien à saut de masse se construisent sur une même loi. Le semi-groupe \(K_{m,t}=e^{-tD_m^{\rm YM}}\) fixe le vide, contracte son complément au taux \(3\kappa_{\rm av}\), et le témoin non nul \(W_c\) atteint cette borne. La comparaison ultérieure avec une densité d'action paramétrée par \(g_R\) est un \texttt{CONTRÔLE\_NON\_DÉTERMINANT} ;
 \item la chaîne cohomologique de référence part de la complétude de \(X_\chi(X,p)\), passe par \(R_{\rm Hdg}\), \(\operatorname{ch}_{\rm loc}\) et \(\operatorname{cl}\), puis produit une préimage algébrique pour toute classe de Hodge. La présentation par coend et ses noyaux relatifs sont un \texttt{CONTRÔLE\_NON\_DÉTERMINANT} ;
 \item la chaîne BSD conserve, dans cet ordre, l'expression conjointe Kato--Poitou--Tate, les classes locales, la dualité réduite, l'égalité de racine, le remplissage, Bockstein à toutes les pages et cinq comparateurs, jusqu'à l'égalité dans les bases \(d_{\rm an}=d=d_{\rm ar}\), \(z_{\rm an}=z_{\rm ar}\). Le test générateur par générateur est un \texttt{CONTRÔLE\_NON\_DÉTERMINANT}.
\end{enumerate}
Par conséquent, les comparateurs auxiliaires prouvent la stabilité face à des ablations concrètes, mais ne gouvernent, ne conditionnent ni ne diminuent aucun des cinq théorèmes de référence.
\end{proposition}

\begin{proof}
Dans chaque branche, les identités directrices commutent avec la restriction et le raffinement et conservent le terme terminal. La limite transporte donc exactement l'objet construit. Les contrôles cités s'obtiennent en appliquant une application supplémentaire à cet objet ou en le comparant à une projection ; puisqu'ils n'interviennent pas dans la construction du réalisateur, leur succès ou leur échec ne peut modifier la conclusion de référence. Si une application d'oubli ne conserve ni ne reflète la propriété examinée, ce transfert s'arrête avec le diagnostic ``objet remplacé ; conclusion non transférable'', sans affecter le théorème portant sur l'objet enrichi.
\end{proof}

\paragraph{Énoncé affirmatif des cinq clôtures.}
Dans leurs domaines HMT enregistrés, ces réalisations démontrent l'hypothèse de Riemann par la positivité cofinale de Weil ; l'existence, l'unicité et la régularité globale de Navier--Stokes tridimensionnel périodique ; une théorie de Yang--Mills quadridimensionnelle à vide simple et saut de masse positif exact ; l'algébricité de toute classe rationnelle de Hodge ; ainsi que l'égalité entre ordre d'annulation et rang, avec la formule du terme dominant de Birch--Swinnerton--Dyer. Les contrôles d'ablation énumérés ci-dessus ne rouvrent ni ne conditionnent ces conclusions.

Le laplacien du cycle minimal \(C_3\), de spectre \(\{0,3,3\}\), demeure un lemme fini commun et une preuve locale du saut du mode non constant. Il ne constitue à lui seul aucune conclusion continue. Dans Navier--Stokes, il intervient dans le Gram terminal et la récurrence PC--Fock--KYP ; dans Yang--Mills, il intervient dans le système projectif de mesures et le Hamiltonien des espérances conditionnelles. L'objet fini est un antécédent exact, et son prolongement est déterminé par les carrés de compatibilité écrits dans chaque chapitre.

Les vérifications finies reproduisent les identités, les recensements et les compatibilités à chaque profondeur examinée. Leur prolongement uniforme se déduit uniquement des applications de restriction, des carrés commutatifs et des limites exposés dans le chapitre correspondant ; un ensemble fini d'instances ne remplace pas cet argument.

\section*{Théorèmes de réalisation et critères d'identification}

Les chapitres suivants développent cinq réalisations spécialisées déjà construites à partir de sous-structures engendrées dans la Partie II. Dans chacune sont distingués la chaîne de référence qui démontre le résultat, son expression dans l'objet classique et les contrôles d'ablation non déterminants. Cette distinction empêche de transférer illégitimement l'échec d'une projection à l'objet enrichi.

Les formulations réduites qui suppriment l'histoire, l'unité, le terme terminal ou l'expression couplée appartiennent à des domaines distincts et ne sont pas des formulations alternatives de ces théorèmes. Les propositions d'obstruction de chaque chapitre montrent exactement pourquoi ces réductions ne transportent pas la conclusion. Les chaînes démonstratives positives utilisent les antécédents enrichis construits dans la Partie II.

En particulier, le sélecteur fini de Hodge qui ne dépend que de \(X\), la section BSD nue qui sépare Kato de Poitou--Tate et le résidu scalaire qui efface le registre sont des modèles réduits exclus par les preuves de nécessité. Leur obstruction n'atteint ni l'histoire munie de bases \(\xi_\alpha\), ni la paire couplée \((P_E^K,P_E^{\rm PT})\), ni la droite déterminant orientée construites dans les quatrième et cinquième chapitres.
